\documentclass[10pt,a4paper]{amsart}
\usepackage[margin=19mm]{geometry}
\usepackage{amsmath,amssymb,mathtools}
\usepackage[scr=rsfs,cal=euler]{mathalfa}
\usepackage{xfrac}
\usepackage[unicode,hidelinks,bookmarks=false]{hyperref}
\newcommand{\ie}{i.e.\ }
\newcommand{\cf}{cf.\ }
\newcommand{\resp}{respectively\ }
\newcommand{\Subsection}{\subsection}
\providecommand{\nobreakdash}{\nobreak\hbox{-}}
\setlength{\emergencystretch}{2em}
\multlinegap0pt
\usepackage{amssymb}
\usepackage{mathtools}
\usepackage[shortlabels]{enumitem}
\newtheorem{theorem}{Theorem}
\newtheorem{lemma}{Lemma}
\newcommand{\E}{\mathbb E}
\newcommand{\TV}{\operatorname{TV}}
\newcommand{\calG}{\mathcal G}
\newcommand{\calP}{\mathcal P}
\newcommand{\calQ}{\mathcal Q}
\newcommand{\calS}{\mathcal S}
\title{Star observations in bounded-degree graphs}
\author{Balazs Szegedy}
\date{Source: arXiv:2605.10785v2 (CC BY 4.0)}
\begin{document}
\maketitle

\noindent\emph{Source and licence.} Star observations in bounded-degree graphs. Version arXiv:2605.10785v2 (CC BY 4.0), distributed by arXiv under the Creative Commons Attribution 4.0 International licence, http://creativecommons.org/licenses/by/4.0/, as recorded in the arXiv licence metadata for this exact version and shown on the abstract page https://arxiv.org/abs/2605.10785v2. Full licence notice: \texttt{licenses/arxiv-2605.10785-CC-BY-4.0.txt}.\par\smallskip

\noindent\textit{Editorial scope.} Counted unit: the article's theorem thm:HLS (source0.tex lines 155--161). The article's complete original proof of it is delivered with it (source0.tex lines 609--711, one \texttt{proof} environment of 103 lines, which the article places away from the statement). No mathematics is written, repaired or re-derived: the statement and the proof are the article's own lines, in the article's own notation, and no retained line is substituted or re-flowed.  Retained uncounted as the source-local context the proof needs: lines 184--190; lines 210--216; lines 277--283; lines 325--335; lines 467--478; lines 509--528; lines 563--573.  Nothing was omitted from any retained span, so the package carries no omission record.  No retained line contains a citation, so the wrapper carries no bibliography and no bibliography entry.  No retained line contains a figure environment, an image-inclusion command or drawing code, so no image is delivered and the package declares no asset dependency.  Editorial formatting only, in addition: an AMS \texttt{amsart} preamble that restates the article's own declarations and macros used by the retained lines, namely the environments \texttt{theorem}, \texttt{lemma} on separate counters, together with the standard packages those lines need.  The retained environments print in this document as Theorem 1, Lemma 1 in order of appearance, whereas the article prints the same items with the numbering of its own full document.  The complete native source of the article as distributed in the arXiv e-print for this exact version is bundled unchanged as \texttt{sources/200320/source0.tex}.
\begin{lemma}[Continuity of star statistics]\label{lem:star-continuity}
If $G_n$ converges locally-globally to a graphing $G$, then, for every $k\geq 1$,
\[
 \calS_k(G_n)\longrightarrow \calS_k(G)
\]
in Hausdorff distance.
\end{lemma}

\begin{theorem}[Star theorem]\label{thm:main}
Let $(G_n)_{n\geq 1}$ be a sequence of graphings of maximum degree at most $d$.  Then $(G_n)$ is local-global convergent if and only if, for every $k\geq 1$, the closed sets
\[
 \calS_k(G_n)\subseteq\calP(\mathsf S_{d,k})
\]
converge in Hausdorff distance.
\end{theorem}

\begin{lemma}\label{lem:true-neighbourhoods-consistent}
Let $G=(X,\nu,E)$ be a graphing and let $g\colon X\to[k]$ be a coloring such that every radius-$r$ ball has pairwise distinct colors.  Define
\[
 m(x)=\gamma_{r,G,g}(x)\in\calG^*_{d,k,r}.
\]
Then $m$ is consistent at every $x\in X$.
\end{lemma}

\begin{lemma}\label{lem:consistency-from-stars}
For every $d,k,r$ and every $\varepsilon>0$ there exists $\delta>0$ with the following property.  Let $G=(X,\nu,E)$ and $G'=(X',\nu',E')$ be graphings, and let
\[
 m\colon X\to\calG^*_{d,k,r},\qquad m'\colon X'\to\calG^*_{d,k,r}
\]
be measurable maps.  Assume that $m'$ is consistent at every point of $X'$ and that, after identifying the finite set $\calG^*_{d,k,r}$ with a color set,
\[
 d_{\TV}\bigl(\sigma(G,m),\sigma(G',m')\bigr)\leq\delta.
\]
Then, for a random $x\in X$, with probability at least $1-\varepsilon$, the map $m$ is consistent at every vertex of $B_r(x)$.
\end{lemma}

\begin{lemma}\label{lem:isomorphism-from-closed-walks}
Let $G=(X,\nu,E)$ be a graphing, let $m\colon X\to\calG^*_{d,k,r}$ be measurable, and let $s\colon X\to[k]$ be the color of the root in $m(x)$.  Suppose that $m$ is consistent at every vertex of $B_r(x)$ and that
\[
 |W_{q}^{0,G}(x)|=|W_{q}^{0,m(x)}(o)|
 \qquad\text{for every }1\leq q\leq r.
\]
If $1\leq t<r/3$, then
\[
 \gamma_{t-1,m(x)}(o)\cong \gamma_{t-1,G,s}(x)
\]
as rooted colored graphs.
\end{lemma}

\begin{lemma}\label{lem:closed-walk-stability}
For every $d,k,r$ and every $\varepsilon>0$ there exists $\delta>0$ with the following property.  Let $G'=(X',\nu',E')$ be a graphing with a coloring $g'\colon X'\to[k]$ such that every radius-$r$ ball has pairwise distinct colors, and put
\[
 m'(x)=\gamma_{r,G',g'}(x)\in\calG^*_{d,k,r}.
\]
Let $G=(X,\nu,E)$ be another graphing, and let $m\colon X\to\calG^*_{d,k,r}$ be measurable.  If, after identifying $\calG^*_{d,k,r}$ with a finite color set,
\[
 d_{\TV}\bigl(\sigma(G,m),\sigma(G',m')\bigr)\leq\delta,
\]
then, for every $1\leq t\leq r$,

\[
 \E_\nu\left|\,|W_t^{0,m(x)}(o)|-|W_t^{0,G}(x)|\,\right|
 \leq c_t(G')-c_t(G)+\varepsilon,
\]
and consequently
\[
 c_t(G)\leq c_t(G')+\varepsilon.
\]
\end{lemma}

\begin{lemma}\label{lem:closed-walks-equal}
Let $G$ and $G'$ be graphings of maximum degree at most $d$.  Suppose that
\[
 \calS_k(G)=\calS_k(G')
\]
for every $k\geq 1$.  Then
\[
 c_t(G)=c_t(G')
\]
for every $t\geq 1$.
\end{lemma}

\begin{theorem}[Local-global compactness]\label{thm:HLS}
For fixed maximum degree $d$, the space of local-global equivalence classes of graphings of degree at most $d$ is compact and metrizable.  Equivalently, every sequence of degree-$d$ graphings has a local-global convergent subsequence, and every local-global convergent sequence has a graphing limit $G$ such that, for every $k\geq 1,r\geq 0$,
\[
 \calQ_{k,r}(G_n) \longrightarrow \calQ_{k,r}(G)
\]
in Hausdorff distance.
\end{theorem}

\begin{proof}[Proof of Theorem~\ref{thm:main}]
As noted before, local-global convergence implies convergence of the sets of colored star distributions $\calS_k$ by Lemma~\ref{lem:star-continuity}.

Conversely, assume that $\calS_k(G_n)$ converges for every $k\geq 1$.  We prove that $(G_n)$ is local-global convergent.  By the compactness theorem recalled in Theorem~\ref{thm:HLS}, it is enough to show that any two local-global subsequential limits are local-global equivalent.

Let $G$ and $G'$ be two such subsequential limits.  Since the star data converge along the whole sequence and, by Lemma~\ref{lem:star-continuity}, converge along the two subsequences to the star sets of their local-global limits, we have
\[
 \calS_k(G)=\calS_k(G')
 \qquad\text{for every }k\geq 1.
\]
By Lemma~\ref{lem:closed-walks-equal},
\[
 c_t(G)=c_t(G')
 \qquad\text{for every }t\geq 1.
\]

We prove $\calQ_{k,R}(G')\subseteq\calQ_{k,R}(G)$ for arbitrary $k,R$; the reverse containment follows symmetrically.

Fix $k\geq 1$, $R\geq 1$, and a measurable coloring
\[
 g'\colon X'\to[k]
\]
of $G'$.  Choose an integer $r$ so large that $R+1<r/3$.  Let
\[
 s_r\colon X'\to[K_{d,r}]
\]
be a coloring that separates all radius-$r$ balls in $G'$, and form the product coloring
\[
 \widetilde g'=(g',s_r)\colon X'\to[k]\times[K_{d,r}].
\]
Let
\[
 m'(x)=\gamma_{r,G',\widetilde g'}(x)
 \in\calG^*_{d,kK_{d,r},r}.
\]
By Lemma~\ref{lem:true-neighbourhoods-consistent}, $m'$ is consistent everywhere.

Put
\[
 \mathcal M=\calG^*_{d,kK_{d,r},r},\qquad a=|\mathcal M|,
\]
choose a bijection $\mathcal M\leftrightarrow[a]$, and regard $m'$ as an $a$-coloring of $G'$.  Since $\calS_a(G)=\calS_a(G')$, for every $\delta>0$ there is a measurable map
\[
 m\colon X\to\mathcal M
\]
such that
\[
 d_{\TV}\bigl(\sigma(G,m),\sigma(G',m')\bigr)\leq\delta.
\]
Let $\widetilde g\colon X\to[k]\times[K_{d,r}]$ be the root color of $m(x)$, and let $g$ be its first coordinate.

By Lemma~\ref{lem:consistency-from-stars}, for all but $o_\delta(1)$ of the points $x\in X$, the map $m$ is consistent on all of $B_r(x)$.  By Lemmas~\ref{lem:closed-walk-stability} and~\ref{lem:closed-walks-equal}, for each $1\leq q\leq r$,
\[
 \E_\nu\left|\,|W_q^{0,m(x)}(o)|-|W_q^{0,G}(x)|\,\right|=o_\delta(1).
\]
Let
\[
 D_q=\{x\,\bigm|\, |W_q^{0,m(x)}(o)|\neq |W_q^{0,G}(x)|\}.
\]
The discrepancy is integer-valued, hence $\nu(D_q)\leq \E  \bigl|\,|W_q^{0,m(x)}(o)|-|W_q^{0,G}(x)|\,\bigr|=o_\delta(1)$.  Taking the union over the finitely many $q=1,\ldots,r$, we obtain a set of measure $o_\delta(1)$ outside which
\[
 |W_q^{0,m(x)}(o)|=|W_q^{0,G}(x)|
 \qquad\text{for all }1\leq q\leq r,
\]
and $m$ is consistent on $B_r(x)$.  On this good set, Lemma~\ref{lem:isomorphism-from-closed-walks}, applied with $t=R+1$, gives
\[
 \gamma_{R,m(x)}(o)
 \cong
 \gamma_{R,G,\widetilde g}(x).
\]

The root marginal of $\sigma(G,m)$ is the law of $m(x)$, and the root marginal of $\sigma(G',m')$ is the law of $m'(x')$.  Therefore these two laws are $o_\delta(1)$-close.  Applying the deterministic truncation map
\[
 \mathcal M\to \calG_{d,kK_{d,r},R},\qquad M\mapsto \gamma_{R,M}(o),
\]
we see that the distribution of $\gamma_{R,m(x)}(o)$ is $o_\delta(1)$-close to the distribution of
\[
 \gamma_{R,G',\widetilde g'}(x').
\]
Thus
\[
 \tau_R(G,\widetilde g)
\]
can be made arbitrarily close to
\[
 \tau_R(G',\widetilde g').
\]
Finally let
\[
 \Pi\colon \calP(\calG_{d,kK_{d,r},R})\to \calP(\calG_{d,k,R})
\]
be the affine continuous map induced by forgetting the auxiliary separating color.  Applying $\Pi$, we conclude that
\[
 \tau_R(G,g)
\]
can be made arbitrarily close to
\[
 \tau_R(G',g').
\]
Thus $\tau_R(G',g')\in\calQ_{k,R}(G)$.  Since $g'$ was arbitrary, $\calQ_{k,R}(G')\subseteq\calQ_{k,R}(G)$.

Interchanging the roles of $G$ and $G'$ gives the reverse containment.  Hence the two graphings are local-global equivalent.  Therefore all local-global subsequential limits of $(G_n)$ are equivalent, and $(G_n)$ is local-global convergent.
\end{proof}

\end{document}
